\chapter{Experimental Design}
\label{chap:experimental-design}

\section{Testing with a Known, Controllable Chirp Source}
\label{sec:exp-motivation}

% TODO: Why a glass window (or stack of them) at an angle is a good
% controllable test source for spatial chirp -- known, adjustable via
% window count/angle, independent of the laser's own chirp.

\section{Glass-Induced Spatial Chirp}
\label{sec:exp-glass-chirp}

% TODO: n = n(lambda) at incidence angle theta -> wavelength-dependent
% refraction -> angular chirp -> (via chapter 1's focus conversion) spatial
% chirp at the spectrometer. See presentation slide 30.

\section{Predicted Spatial Dispersion}
\label{sec:exp-predicted-dispersion}

% TODO: Derivation of zeta_predicted from window thickness/count, angle,
% and glass dispersion (Sellmeier or similar) -- the values quoted in
% chapter 7 (e.g. zeta_predicted = 81.9 +/- 2.8 nm/nm for 1 window) should
% trace back to this derivation.

\section{Physical Setup}
\label{sec:exp-physical-setup}

% TODO: Series of glass windows placed upstream of the imaging
% spectrometer on the bench. See presentation slides 31-32 (schematic +
% bench photo).

\section{Measurement Procedure}
\label{sec:exp-procedure}

% TODO: How each configuration (0/1/2/3 windows) was measured -- number of
% shots, baseline subtraction procedure, what "baseline subtracted" means
% for the 1/2/3-window results in chapter 7.
